\documentclass[11pt,a4paper]{article}
\usepackage[margin=24mm]{geometry}
\usepackage[T1]{fontenc}
\usepackage{lmodern,amsmath,amssymb,amsthm,mathtools}
\usepackage{booktabs,longtable,array,multicol,enumitem,microtype,xcolor}
\usepackage[colorlinks=true,linkcolor=blue!45!black,urlcolor=blue!45!black,citecolor=blue!45!black]{hyperref}
\usepackage{fancyhdr}
\pagestyle{fancy}\fancyhf{}\fancyhead[L]{\small U208: a conditional tape-context fold}\fancyhead[R]{\small 29 October 2026}\fancyfoot[C]{\thepage}
\setlength{\headheight}{14pt}
\setlength{\emergencystretch}{2em}
\newtheorem{theorem}{Theorem}[section]
\newtheorem{lemma}[theorem]{Lemma}
\newtheorem{proposition}[theorem]{Proposition}
\newtheorem{corollary}[theorem]{Corollary}
\newcommand{\PA}{\mathrm{PA}}
\newcommand{\Con}{\operatorname{Con}}
\newcommand{\Halt}{\operatorname{Halt}}
\newcommand{\Uold}{U_{209}}
\newcommand{\Unew}{U_{208}}
\title{\textbf{A 208-state consistency machine}\\[3mm]\large Comparison of the U209 proposals and a new conditional tape-context reduction}
\author{Research prepared for Vladimir Reshetnikov}
\date{29 October 2026}
\begin{document}
\maketitle
\begin{abstract}
We construct an explicit binary, single-tape Turing machine with
\textbf{208 working states}, including its blank initializer, by deleting
one state from the supplied U209 candidates. The reduction changes only
two retained transition rows. Its new ingredient is a conditional tape
invariant: when an existing helper reads zero, its predecessor cell is
one, whereas the deleted state has the two-cell predecessor word $10$.
Four exact local identities, of lengths $7,1,3,1$, permit the helper to
serve both roles. A five-state inductive invariant and 414 ordinary row
identities prove halting equivalence on every common initial tape from
state~0. For blank computations the equivalence is formalizable in PA;
composition with the retained unary-theory calibration gives
$\PA\vdash\neg\Halt(\Unew,\mathrm{blank})\leftrightarrow\Con(\PA)$.

The six supplied submissions represent only two machines up to rooted
labeled graph isomorphism. Both admit the new reduction, yielding two
nonisomorphic 208-state variants. The primary candidate passes a fresh
replay of all 5,084 symbolic primitive contracts, with all 416 rows
covered, and retains the complete 48,140-step startup trajectory.
The archive contains mathematical proofs and executable certificates;
it does not contain a newly exported end-to-end proof-assistant
certificate for the consistency biconditional.
\end{abstract}

\section{Result and machine conventions}
The primary artifact is \path{machine/u208.tm}. Its working states are
$0,\ldots,207$ and its separate halt label is 208. The tape alphabet is
$\{0,1\}$, the tape is indexed by $\mathbb Z$, and the initial
configuration has state~0, head coordinate~0, and all cells zero.
Every working-state/symbol pair has exactly one transition and every
transition moves left or right. A transition into the halt label performs
its write and move before stopping. The initializer is included in the
208-state count.

A table row \texttt{q b w D r} means: in state $q$, reading $b$, write
$w$, move one cell in direction $D$, and enter state $r$.
The header \texttt{\#! start 0} is metadata, not a transition.
The proofs use sparse \emph{native} state names, in which 211 is a
working state and 254 is halt. \path{machine/state_map.json} gives the
complete bijection from retained native names to dense labels.

\begin{center}
\begin{tabular}{@{}lr@{}}
\toprule
Property & Primary candidate\\
\midrule
Working states / transition rows & $208/416$\\
Changed retained rows / deleted rows, relative to Joint U209 & $2/2$\\
Primitive source occurrences / guarded ordinal control edges & $3,909/5,083$\\
Changed source addresses or primitive operands & $0$\\
Symbolic primitive contracts / covered literal rows & $5,084/416$\\
Maximum local time expansion per old transition & $7$\\
Exact unchanged blank startup length & $48,140$\\
\bottomrule
\end{tabular}
\end{center}

The dense table has SHA-256
\begin{center}\small\ttfamily
a2fee9e41311c0d43aaac871ea411be24eada17e6b3633c55f40b437101144e1.
\end{center}
The native table has SHA-256
\begin{center}\small\ttfamily
b348a466a2c316511e0037ea3824de29d4c6ad37338404fbc98a3968fdb25135.
\end{center}
The deterministic builder, the independent local checker, and the full
table in Appendix~\ref{app:table} identify the candidate independently of
the prose. Its local reduction is proved without a consistency
assumption. Actual nonhalting is equivalent to consistency of PA;
neither unconditional consistency nor global state minimality is claimed.

The scope of novelty is explicit: this is a further reduction of the
six supplied U209 submissions at the inspected repository snapshot.
The sources did not contain the new $172\mapsto211$ fold. The inspection
is not an exhaustive search of all unpublished machines or a priority
claim against the entire literature.

\tableofcontents
\clearpage
\section{Comparison of the supplied proposals}
\label{sec:comparison}
\subsection{Frozen scope and exact identities}
The incoming directory was retrieved at ProveIt commit
\begin{center}\small\ttfamily
343da82d7384d10ec9fd1c17361fc923b2d1d117.
\end{center}
It contained five executable U209 packages, a sixth U209 ZIP containing
only a PDF, and the unary-theory archive~\cite{unary,plain,compat,folds,composed,joint,jointpdf}.
The download of every item was checked against its Git blob identifier.
\path{reference/source_inventory.json} retains the exact URLs, names,
sizes, and blob identifiers; the comparison snapshots also record SHA-256.

\begin{center}
\small
\begin{tabular}{@{}llcl@{}}
\toprule
Submission & Dense SHA-256 prefix & Class & Source depth\\
\midrule
Plain U209 & \texttt{65ccb606c8e9} & A & 42\\
Compatibility & \texttt{bda91bd40880} & B & 41\\
Compatible Folds & \texttt{5ded628d692f} & A & 42\\
Composed & \texttt{3a69998d818b} & A & 42\\
Joint, full package & \texttt{5ded628d692f} & A & 42\\
Joint, PDF-only package & \texttt{65ccb606c8e9} & A & 42\\
\bottomrule
\end{tabular}
\end{center}

Here ``class'' means rooted deterministic transition-graph isomorphism
preserving the initial state, halt, read symbols, write symbols, and
directions. Compatible Folds and full Joint have byte-identical native
and dense tables. Plain and Composed differ from them only by state
renaming. The four class-A executable packages also have byte-identical
source occurrence maps. The PDF-only Joint appendix was extracted
independently: all 418 printed transition rows agree with Plain. It
therefore identifies Plain's table, even though its title resembles
the full Joint package.

There are consequently \emph{two}, rather than five or six, distinct
machines in the comparison. A different archive, table hash, or rotation
recipe need not describe a different machine.

\subsection{Independent isomorphism and source checks}
A rooted read-labeled graph map is forced by following zero and one
edges from state~0. An independent breadth-first procedure propagates
that map, checking writes, directions, injectivity, and the distinguished
halt label. Every working state is graph-reachable in every table, so
a consistent map covers the whole machine. This is syntactic graph
reachability; it does not assert that the actual blank computation
visits every state.

In native labels, the nontrivial permutations are particularly small:
\[
\begin{array}{ll}
\text{Plain to Joint or Compatible Folds:}&
 (217\ 129\ 214)(215\ 204\ 216),\\
\text{Plain to Composed:}&(215\ 129\ 204\ 214\ 216).
\end{array}
\]
All other labels are fixed, including initial state~0 and halt~254.
The checker verifies the conjugacy of every row. For Plain versus
Compatibility, the forced map has an explicit conflicting assignment;
the witness is retained in \path{comparison/independent_comparison.json}.
Its edge words are paths in a labeled graph, not claimed tape histories.

The two source layouts differ in exactly 720 addresses. An independent
interval-based resolver verifies the same ordered 3,909 opcode/operand
pairs and all 5,083 guarded ordinal control edges. Both have maximum
consecutive-zero run six and 197 distinct source subtrees. Compatibility
reduces maximum address depth from 42 to 41. A shallower source tree alone
does not establish a better asymptotic runtime for the complete proof search.

\subsection{What is shared, and what each package contributes}
The U209 machines descend from one common U212 construction. Their
joint layouts allow deletion of native states 180, 181, and 218, which
had previously been removed separately in three U211 proposals.
Class A uses representatives 163, 220, and 157. Compatibility instead
uses 216 for state 180, with 220 and 157 serving the other two folds.
All packages retain the arithmetic-entry repair at native row $171/0$,
redirecting the backward entry from 7 to 205. That repair is necessary
because a source decoder also serves an arithmetic role.

The common reduction method exploits bits written immediately before
a decoder entry. Opposite predecessor bits allow two control roles to
share a state when an existing helper recovers the role from the tape.
The source reassociations ensure that the required tags coexist without
breaking the primitive control-flow graph or other interpreter entries.
Different deleted labels therefore do not automatically give additive
savings: the helper and address requirements can conflict.

\paragraph{Plain.}
This package gives a detailed ledger of 30 closed source rotations from
U212, explicit local quotient cases, and careful treatment of the
arithmetic entry. It distinguishes algebraic intermediate layouts from
valid final interpreters: some intermediate addresses have seven-zero
runs, while the final runtime proof uses the restored six-zero bound.
It also supplies a useful obstruction: simply deleting the apparently
dominated propositional row destroys the intended PA calibration.
A Boolean interpretation using equivalence as implication validates
the remaining calculus while keeping its target false. That is a
reason to retain the row, not an unexplored one-state saving.

\paragraph{Compatibility.}
This is the genuinely different U209 family. Its eight-factor
reassociation changes 720 addresses relative to class A and produces
depth 41. Its separate construction, runtime, and theory audits make
the changed source layout easy to inspect. It provides an important
second search starting point, though its numerical calculus and
symbolic inference engine are shared with the other packages.

\paragraph{Compatible Folds.}
Its staged 212-to-211-to-210-to-209 construction makes the coexistence
of the old three folds explicit. It is a useful presentation of the
same literal machine later delivered in the full Joint package.
The extra exposition is valuable even though it is not an additional
transition-graph class.

\paragraph{Composed.}
The grouped source proof and independent local-fold checker provide
another view of class A. The machine itself is a permutation of Plain's
decoder labels; its savings cannot be added to Plain's savings.

\paragraph{Joint, full package.}
This is the selected baseline for the new construction. Its direct
18-selector layout, compact complete ancestry, independent source-prefix
auditor, and finite C++ executor make it convenient to extend and review.
The selection is about its evidence organization, not a smaller state count
than the other class-A presentations.

\paragraph{Joint, PDF-only package.}
The report is readable and its printed table is complete, but it does
not add an executable verification package. Its table agrees with
Plain, not byte-for-byte with full Joint. This distinction matters when
binding a mathematical argument to a literal transition file.

\subsection{Fresh comparison evidence and remaining common dependencies}
All five executable packages passed their advertised quick replay
commands in separate copies. Those commands check manifests, deterministic
construction, local folds, source bindings, and startup according to their
documented scope. The six-table comparison was independently rerun under
ordinary and optimized Python. The comparison phase did not rerun five
copies of the full symbolic runtime derivation. Instead, the new primary
U208 table received its own complete fresh derivation in
Section~\ref{sec:verification}.

The inherited runtime engines are essentially the same implementation
under label and import changes. Multiple passing receipts should not be
presented as independent logical foundations. The unary-theory
interpretation, native numerical calculus, strong compiler, and protected
startup argument are also shared. The new reduction preserves these
components and has a short separate local proof.

\subsection{Why the new fold was outside the old local search}
The retained Joint search tested direct aliases through 15 steps and
one-row folds using opposite unconditional predecessor tags through
17 steps. No result was found in those scopes. Native state 211 is
entered from both directions and does not possess a uniform predecessor
tag. Its useful property is conditional on reading zero; state 172
needs two predecessor cells. Extending the search to finite-history
contexts exposes this interface. The successful proof needs only the
five-state invariant and the four finite identities below, and does
not depend on search bounds or random screening.

\section{A conditional tape-context fold from 209 to 208 states}
\label{sec:local208}

The reduction in this section acts directly on the supplied Joint U209
transition table. It preserves the tape alphabet, the blank initializer, and
the entire numerical source program. Its additional information is a property
of the recent computation: two states with different zero branches can share
one physical state because their reachable tape contexts distinguish those
branches. The required property of one state is conditional on its scanned
symbol. The opposite-predecessor-bit interfaces considered in the supplied
searches do not include this case.

All state labels in this section are \emph{native} labels. In particular, 211
is a working state and 254 is the separate halt label. Dense renumbering is
performed only after the reduction. Let $\Uold$ denote the Joint parent and
$\Unew$ the reduced machine.

\subsection{The complete modification}
\label{subsec:local208-patch}

Delete state 172 and redirect its sole incoming transition to state 211.
Change the zero transition of state 211 as follows:
\begin{center}
\begin{tabular}{llll}
\toprule
Row & Parent action & Reduced action & Purpose\\
\midrule
$116/0$ & $0\,R\,172$ & $0\,R\,211$ & Redirect the deleted state\\
$211/0$ & $0\,R\,157$ & $0\,L\,217$ & Inspect the preceding tape context\\
\bottomrule
\end{tabular}
\end{center}
The removed rows are $172/0:0\,L\,67$ and $172/1:1\,R\,184$.
Every other retained row is unchanged. Thus 209 working states and 418
transition rows become 208 working states and 416 rows. The initial state
remains 0, and no halt transition is added or removed.

For the simulation, define the state map
\[
 f(q)=\begin{cases}211,&q=172,\\q,&q\ne172.\end{cases}
\]
The map identifies two states, but changes neither the head position nor
the tape. Its lack of injectivity is harmless: the proof follows a definite
parent computation and establishes agreement at explicitly defined times.

\subsection{Reachable contexts, proved by induction}
\label{subsec:local208-guards}

A configuration is $(q,h,T)$, where $h\in\mathbb Z$ and
$T:\mathbb Z\to\{0,1\}$. Write $t_j=T(h+j)$ for the tape viewed from
the current head. The following five conditions are imposed only when the
configuration has the indicated state:
\begin{align*}
 I_{116}:&\quad t_{-1}=1, &
 I_{172}:&\quad t_{-2}=1\ \land\ t_{-1}=0,\\
 I_{217}:&\quad t_{-1}=1, &
 I_{163}:&\quad t_{-1}=1\ \lor\ t_{-2}=1,\\
 I_{211}:&\quad t_0=1\ \lor\ t_{-1}=1.
\end{align*}
Let $I(q,h,T)$ be the conjunction of the five corresponding implications
from the state equality. In particular, $I_{211}$ says that a scanned zero
at state 211 forces the cell immediately to its left to be one. It does
\emph{not} assert that every entry to 211 has a left-hand one.

The entire incoming-edge census needed for these conditions is small:
\begin{center}
\begin{tabular}{cl}
\toprule
Target & Every parent transition entering the target\\
\midrule
172 & $116/0:0\,R\,172$\\
116 & $215/1:1\,R\,116$\\
211 & $163/0:0\,L\,211$, \quad $163/1:1\,R\,211$\\
163 & $132/1:1\,R\,163$, \quad $199/1:1\,R\,163$,
      \quad $217/0:0\,R\,163$\\
217 & $216/1:1\,R\,217$\\
\bottomrule
\end{tabular}
\end{center}
These are inventories of all physical transition rows, including entries
arising outside ordinary source decoding. The argument therefore does not
depend on a sampled run or an assumption that additional entrances are
unreachable.

\begin{lemma}[Five-state invariant]
\label{lem:local208-invariant}
Every configuration of $\Uold$ reached from state 0 satisfies $I$, on
every initial binary tape.
\end{lemma}
\begin{proof}
The initial configuration satisfies $I$ because 0 is none of the five
listed states. Assume $I$ holds before a transition. Conditions whose
target state is not reached impose no obligation afterward, so it suffices
to inspect the displayed census.

An entry to 116 or 217 writes one and moves right, establishing its
left-hand-one condition immediately. An entry to 172 comes from 116/0.
It writes zero immediately left of the new head; the previous left-hand
one supplied by $I_{116}$ becomes the cell two positions left of that head.
This proves $I_{172}$.

An entry to 163 from 132/1 or 199/1 writes one immediately left of the new
head. Its remaining entrance, 217/0, writes a left-hand zero but carries
the one supplied by $I_{217}$ to the position two cells left. Each case
establishes the disjunction $I_{163}$.

Finally, entry to 211 through 163/1 writes one immediately left of the
new head. Entry through 163/0 moves left instead. The new current cell is
the old $t_{-1}$, and the new left cell is the old $t_{-2}$. Their disjunction
is exactly $I_{163}$, so $I_{211}$ follows in this case as well. This proves
preservation of all five conditions and completes the induction.
\end{proof}

The invariant explains the available distinction. State 172 necessarily
arrives immediately to the right of the word $10$. State 211, whenever it
reads zero, instead has a one immediately to its left. Recovering this
information needs at most three preceding parent transitions. Nevertheless,
the inductive formulation is preferable for composition: it applies uniformly
at every time, including the beginning of the run, without a separate
assumption that a sufficiently long history already exists.

\subsection{Four exact finite identities}
\label{subsec:local208-blocks}

At a simulation boundary the reduced machine has state $f(q)$, head $h$,
and the same tape $T$ as the parent. Only parent states 172 and 211 require
special treatment. The four contracts below specify the complete effects
of their replacement blocks. Head displacements are relative to $h$.
\begin{center}
\begin{tabular}{ccccc}
\toprule
Parent state & Read & Additional tape condition & Steps & Final state, head\\
\midrule
172 & 0 & $t_{-2}=1,\ t_{-1}=0$ & 7 & $67,\ -1$\\
172 & 1 & None & 1 & $184,\ +1$\\
211 & 0 & $t_{-1}=1$ & 3 & $157,\ +1$\\
211 & 1 & None & 1 & $184,\ +1$\\
\bottomrule
\end{tabular}
\end{center}
Every contract preserves the entire tape. The additional conditions for
the zero cases follow from Lemma~\ref{lem:local208-invariant}. The one
cases need only the stated scanned bit, even though other invariant
conditions may also hold.

\begin{lemma}[Local replacement]
\label{lem:local208-replacement}
Suppose a nonhalting parent configuration $C=(q,h,T)$ satisfies $I$,
and its successor is $C'=(q',h',T')$. From $(f(q),h,T)$ the reduced machine
reaches $(f(q'),h',T')$ in one, three, or seven transitions. All earlier
configurations of this block are nonhalting.
\end{lemma}
\begin{proof}
If $q\notin\{172,211\}$, the parent and reduced rows agree after applying
$f$ to the target. This includes row 116/0, whose target is precisely the
one redirected by $f$. There are 414 such ordinary rows.

For either exceptional state reading one, the reduced machine uses the
unchanged row $211/1:1\,R\,184$. Both parent states have that same row,
so one transition suffices.

For parent state 211 reading zero, its invariant supplies $t_{-1}=1$.
The new row moves left into 217; row $217/1:1\,R\,215$ moves back to
the original zero; row $215/0:0\,R\,157$ gives the endpoint of the
parent's $211/0$ transition. The three transitions read $0,1,0$ and write
those same bits.

For parent state 172 reading zero, the invariant supplies
$(t_{-2},t_{-1},t_0)=(1,0,0)$. The full reduced computation is:
\begin{center}
\begin{tabular}{rrrrrr}
\toprule
Step & State before & Head before & Read/write & State after & Head after\\
\midrule
1 & 211 & $0$  & $0/0$ & 217 & $-1$\\
2 & 217 & $-1$ & $0/0$ & 163 & $0$\\
3 & 163 & $0$  & $0/0$ & 211 & $-1$\\
4 & 211 & $-1$ & $0/0$ & 217 & $-2$\\
5 & 217 & $-2$ & $1/1$ & 215 & $-1$\\
6 & 215 & $-1$ & $0/0$ & 157 & $0$\\
7 & 157 & $0$  & $0/0$ & 67  & $-1$\\
\bottomrule
\end{tabular}
\end{center}
Its endpoint is exactly the parent transition $172/0:0\,L\,67$.
Every read is of one of the three specified cells and every write preserves
that cell. The three-step case has the same property for its two cells.
No unspecified exterior cell is inspected, and no exceptional block visits
halt label 254. This proves all cases.
\end{proof}

The seven-step block visits the merged state 211 twice. There is no
circular proof here: all seven transitions have been evaluated explicitly.
In particular, the argument does not assume that arbitrary internal
configurations of the reduced computation satisfy the parent invariant.
The invariant is required only at parent simulation boundaries.

\subsection{Global equivalence and an exact clock}
\label{subsec:local208-global}

For a parent computation, let $N_{172,0}(n)$ count occurrences of state 172
reading zero among its first $n$ transitions, and define $N_{211,0}(n)$
similarly. Counts are taken along the parent run, not the internal states
of the reduced blocks. Define
\begin{equation}
\label{eq:local208-clock}
 \tau(n)=n+6N_{172,0}(n)+2N_{211,0}(n).
\end{equation}
The excesses six and two are the extra transitions used by the seven-step
and three-step replacements.

\begin{theorem}[Uniform halting equivalence]
\label{thm:local208-equivalence}
For every binary tape $T$ and initial head position $h$, $\Uold$ and
$\Unew$, both started in state 0 on $(h,T)$, either both halt or both
run forever. At every parent time $n$ up to and including its first halt,
the reduced configuration at time $\tau(n)$ has the same head and tape,
and its state is the image under $f$ of the parent state. Furthermore,
\[
 n\le\tau(n)\le7n.
\]
If the parent first halts at time $H$, the reduced machine first halts
at exactly $\tau(H)$.
\end{theorem}
\begin{proof}
The configurations agree initially because $f(0)=0$. Induct on the number
of parent transitions. Lemma~\ref{lem:local208-invariant} supplies the
required invariant at each parent configuration, and
Lemma~\ref{lem:local208-replacement} extends the agreement by the
appropriate positive block. The block lengths give
\eqref{eq:local208-clock}. Each length is between one and seven, yielding
the bounds. No block contains an earlier halt, so a parent first halt
produces a reduced first halt exactly at the corresponding boundary.

If the parent never halts, the same induction supplies nonhalting reduced
blocks at all boundaries. Their strictly positive lengths make
$\tau(n)$ unbounded, and together the blocks cover every reduced time.
The reduced machine therefore cannot halt between boundaries or after
all of them. This also proves the converse implication of the halting
equivalence. Determinism ensures that these blocks are prefixes of the
actual reduced run, rather than merely possible executions.
\end{proof}

Thus the result concerns the whole computation, not only a finite prefix.
It also supplies quantitative finite statements: a parent halt within $n$
steps implies a reduced halt within $7n$ steps; a reduced halt within $n$
steps implies a parent halt within $n$ steps. No consistency assumption
appears anywhere in the reduction.

\subsection{Arithmetic scope and the other U209 family}
\label{subsec:local208-scope}

The arbitrary-tape theorem above is a mathematical statement about tape
functions. For the required blank input, configurations and finite runs
have the usual arithmetic encodings. The five invariant checks, the four
finite identities, and the induction on coded computation length are
formalizable in $\PA$. The explicit bound $7n$ removes any need for an
unbounded search to locate a corresponding finite prefix. Consequently,
this construction supplies the PA-formalizable equivalence
\[
 \PA\vdash
 \bigl(\Halt(\Uold,\mathrm{blank})
       \leftrightarrow\Halt(\Unew,\mathrm{blank})\bigr).
\]
Composing it with a justified parent consistency calibration transfers
that calibration to $\Unew$. The reduction neither establishes
$\Con(\PA)$ nor replaces the separate proof-search and theory
interpretation obligations. The delivered mathematical proof and finite
certificates do not constitute a newly exported proof-assistant certificate
of the complete consistency endpoint.

The distinct Compatibility U209 family admits the same reduction:
delete 172 into 211 and replace row $211/0$ by $0\,L\,214$.
Its local roles correspond as follows:
\begin{center}
\begin{tabular}{lcc}
\toprule
Role & Joint parent & Compatibility parent\\
\midrule
Left-context selector & 217 & 214\\
State feeding the merged-state entries & 163 & 216\\
Right branch of the selector & 215 & 163\\
Predecessor writing one before the selector & 216 & 217\\
\bottomrule
\end{tabular}
\end{center}
States 116, 172, 211, 157, 67, and 184 retain their roles. The corresponding
complete censuses are: 172 from 116/0; 116 from 163/1; 211 from 216/0
and 216/1; 216 from 132/1, 199/1, and 214/0; and 214 from 217/1.
Their directions and written bits are exactly those used in the invariant
proof. The seven-step state sequence becomes
\[
 211\to214\to216\to211\to214\to163\to157\to67.
\]
The same tape conditions, head positions, block lengths, and clock formula
therefore apply. This produces a second 208-state table from the other
supplied machine family, with no additional source transformation.

\section{The retained consistency calibration}
\label{sec:calibration}

The new reduction preserves the entire source proof search. Its logical
endpoint therefore comes from composing the preceding literal-machine
equivalence with the retained calibration of the selected Joint U209
parent~\cite{joint}. This section identifies the exact mathematical
dependencies and the scope of their verification. The fixed K215 theorem
about ZFC~\cite{k215pa} is background and a source of reusable interfaces;
it is not applied to the changed transition table by renaming its conclusion.

\subsection{Unary power covering}

The proposed theory $\mathsf U$ uses equality-free classical first-order
logic over a nonempty domain, with one binary relation $E$, written as
membership. Put
\[
 \operatorname{Sub}(x,a):=\forall t\,(E(t,x)\to E(t,a)).
\]
For each formula $\varphi$ whose free variables are contained in
$\{x\}$, its sole nonlogical schema is
\begin{equation}
\label{eq:unary-calibration}
\forall a\,\exists b\left[
 \forall x\,(E(x,b)\to\varphi(x))\ \land\
 \forall x\,(\operatorname{Sub}(x,a)\to
                       (\varphi(x)\to E(x,b)))\right].
\end{equation}
All bound variables are freshened before insertion. In particular, the
source $a$ cannot occur freely in $\varphi$. Arbitrarily many bound
variables and unrestricted quantifier complexity are permitted. The cover
contains all selected represented subsets of $a$, and may contain further
selected objects. It need not be an exact filter of the subsets of $a$.
The tempting exactification $\varphi(x)\mapsto
\operatorname{Sub}(x,a)\land\varphi(x)$ is illegal in this unary schema.

The supplied report~\cite{unary} proves interpretations in both directions
between $\mathsf U$ and full PA, with uniform finite-proof transformations.
The reviewed lower argument addresses the parameter restriction through a
specific noncircular sequence of constructions, rather than assuming
unrestricted separation.

Define coextensionality $u\simeq v$ by equality of member profiles. An
object is saturated when its membership predicate respects this relation.
Numerical representatives are saturated, transitive objects whose members
are saturated and linearly ordered by membership modulo coextensionality.
A further endpoint condition says that each nonempty intersection of an
existing object's member profile with the number or one of its numerical
initial segments has least and greatest member classes. All these
conditions are fixed first-order formulas. They assert neither external
finiteness nor external well-foundedness.

False filtering supplies an empty object. The local numerical facts give
heredity, a greatest predecessor for each nonempty number, and uniqueness
of existing successors modulo coextensionality. They do not initially
assert that every number has a successor. A unary filter of failures then
proves induction: if failures occur below a counterexample $n$, choose
their least class and use its predecessor; if there are no failures below
$n$, the predecessor of $n$ supplies the contradiction directly. Neither
case needs an already existing successor of $n$.

For an induction predicate $\theta(v,\bar p)$, write
$A_\theta(\bar p)$ for its base and step premises. Parameter
internalization uses the legal unary formula
\[
 Q(v):=\forall\bar p\,
             (A_\theta(\bar p)\to\theta(v,\bar p)).
\]
This yields induction with arbitrary ambient parameters, even when the
predicate is not invariant under coextensionality. Induction next gives
comparison of numerical classes and least numerical witnesses. Fixed
finite tuples can be minimized lexicographically without assuming objects
that already encode those tuples.

The resulting bounded comprehension theorem has explicit invariance and
support hypotheses. Suppose $\chi(z,\bar n)$ is invariant in its numerical
parameters and every selected $z$ is a subset of some supporting object.
If exact comprehension fails, let $\bar r$ be the least failing numerical
tuple. The unary matrix
\[
 \exists\bar n\,
    (\operatorname{LeastFailure}(\bar n)\land\chi(z,\bar n))
\]
is extensionally $\chi(z,\bar r)$. Filtering at a separately obtained
support source gives the forbidden exact object. The parameters are bound
inside the matrix. This is a finite classical proof construction, not an
algorithm that evaluates a least failure or a hidden truth predicate.

This comprehension supplies total successors, powers of numerical
sources, and saturated singleton, pair, and ordered-pair codes. Pair codes
with coordinates in a number $n$ are subsets of an exact power of $n$,
providing supports for numerical graph comprehension. Deterministic
histories are extended using a canonical formula that existentially binds
the old history; an arbitrary graph is never inserted as a forbidden free
matrix parameter. History uniqueness precedes extension, and a separate
numerical bound controls values without assuming monotonicity. Successor
iteration gives addition, iteration of addition gives multiplication, and
the already derived induction gives every translated PA induction instance.

\subsection{The upper interpretation and internal proof transport}

For the reverse direction, interpret $E(x,a)$ as ``bit $x$ of $a$ is
one.'' PA proves that the translated subset condition implies $x\le a$.
For each fixed arithmetic formula $\chi$, PA induction constructs its
binary characteristic code on the interval through $a$. Applying this to
$\operatorname{Sub}(x,a)\land\varphi(x)$ gives an exact filter and hence
the weak cover~\eqref{eq:unary-calibration}. The argument also permits
arbitrary extra matrix parameters. Because it uses PA's induction scheme,
it applies in nonstandard PA models and requires no uniform satisfaction
predicate. PA also proves global irreflexivity of this membership relation.

Once the lower interpretation is fixed, its finite non-induction
obligations have fixed derivations; an arbitrary induction instance is
inserted into the preceding uniform template. Fresh-name selection,
substitution, operation-graph translation, relativization, and finite proof
concatenation are primitive-recursive syntax operations. The upper
translation similarly inserts each fixed translated formula into the
finite-comprehension proof. Induction on formulas and coded derivations
verifies the resulting refutation transformations in PA. Thus the supplied
claim is the internal syntactic equivalence
\begin{equation}
\label{eq:U-PA-calibration}
 \PA\vdash\Con(\mathsf U)\leftrightarrow\Con(\PA),
\end{equation}
not merely standard-model equiconsistency. The package gives the ordinary
mathematical uniformity argument; it does not provide a complete executable
arbitrary-proof translator or a new proof-assistant certificate for this
assertion.

\subsection{What the native machine actually enumerates}

The selected machine preserves the native Cantor-coded calculus. It does
not newly compile the positive dyadic emitter also proposed in the unary
report. Its ordinary machine-facing theory is $\mathsf W+\mathsf{IR}$:
$\mathsf W$ allows arbitrary matrix parameters in the weak-cover schema,
and $\mathsf{IR}$ is $\forall x\,\neg E(x,x)$. Its consistency calibration
is the cycle of PA-verifiable reductions
\[
 \Con(\PA)\ \Longrightarrow\ \Con(\mathsf W+\mathsf{IR})
 \ \Longrightarrow\ \Con(\mathsf U)
 \ \Longrightarrow\ \Con(\PA).
\]
The first arrow uses the bit interpretation, the second the inclusion of
the unary instances, and the third the lower interpretation just reviewed.

The numerical calculus has truth alias $0$, target $1$, seven logical
rows, a dominated propositional row, modus ponens, generalization, and a
gated cover/vacuity bundle. In the retained role assignment, matrix code
and output-variable label are $a$, the source label is $b$, and the tested
label is $c$; the gate requires $c\ne a,b$. Every effective variable label
in a positive matrix code is smaller than that code, while code zero is
a closed truth alias. The output label is therefore fresh.

When $a\ne b$, the generated guard is the subset condition. When $a=b$,
the guard's local binder captures the source occurrence and gives
emptiness under irreflexivity. A cover at an empty source supplies this
exceptional instance, including the small codes zero and one. The gate
also preserves atomic membership vacuity. For completeness, a desired
matrix is padded to an equivalent code $n>1$ and emitted using
$(a,b,c)=(n,0,1)$; the logical proof templates remove the padding and
restore names. Auxiliary equality is handled through two-sided membership
indiscernibility, not by silently assuming unrestricted primitive-equality
vacuity. These are the retained U214/U213 calibration arguments in the
parent's reference chain~\cite{joint}.

\subsection{Proof search and the protected startup}

The source contains 3,909 primitive occurrences. Its axiom command keeps
the theorem-valued output and ten-entry stack interface needed by the
three-command proof compiler. Strong completeness works above an arbitrary
underlying numerical stack after counter and cache reset. Successful
candidate codes are cofinal, and each residual continuation decreases;
retaining a finite candidate cutoff therefore does not lose completeness.
The native alias $P(0,0)=0$ is explicitly handled, including the special
zero case of the counter-setting macro.

Startup needs a separate soundness argument. The shortened initializer
reaches an interior axiom-command suffix, rather than a complete invocation.
Its source bridge starts at ordinal $106$ with logical store
$(1,0,4,0,0,1,0,0,1)$ and reaches ordinal $136$ after 141 source primitives.
The raw cache value $1$ is used as an arithmetic parameter, and the stack
tail $4$ is opaque; neither is assumed theorem-valued.

The suffix builds nine explicit theorem heads above that tail and reaches
scheduler candidate $1$ with counter and cache zero. A ghost explicit
stack length is retained because the pairing alias prevents recovery of
list length from its integer code. The all-natural scheduler-credit proof
assigns conservative weights $-1,0,10$ to transfer, modus ponens, and axiom
commands. Its maximal prefix debt from candidate $1$ is two, leaving a
permanent margin of seven heads. Hence no pop exposes the opaque tail.
This unbounded argument supplies startup soundness; a finite blank trace
alone would not do so.

\subsection{Composed endpoint and proof status}

\begin{corollary}
For the exact delivered 208-state table, with its own initialization and
blank input tape,
\[
 \PA\vdash
 \neg\Halt(\Unew,\mathrm{blank})
       \leftrightarrow\Con(\mathsf U)
       \leftrightarrow\Con(\PA).
\]
\end{corollary}
\begin{proof}
Compose the new literal halting equivalence of
Theorem~\ref{thm:local208-equivalence} with the selected parent's retained
numerical proof-search, compiler, startup, and theory calibration, then
apply~\eqref{eq:U-PA-calibration}.
\end{proof}

No consistency assumption is needed for the biconditional. Asserting
actual nonhalting additionally uses $\Con(\PA)$. The result is an ordinary
PA-formalizable mathematical proof, supported by exact finite checks and
universal symbolic induction certificates. Those certificates are stronger
than bounded register tests, but their checker is not a proof-assistant
kernel. The separate formal K215/Con(ZFC) and Theory K developments
\cite{k215pa,k230} establish different endpoints. This release does not
claim a newly compiled end-to-end Lean, Rocq, or MM0 theorem for U208.

\section{Verification of the delivered literal machine}
\label{sec:verification}
\subsection{Small independent certificate for the new theorem}
The independent \path{code/check_local.py} imports neither the builder
nor proposal code. It reads the pinned parent, independently constructs
the prescribed child, and verifies its identity with both delivered
native and dense tables. It uses explicit exceptions throughout, so its
checks remain active under \texttt{python3 -O}.

The five-state invariant is checked as a finite Boolean implication for
every possible local tape assignment needed by an incoming transition.
The window $[-3,1]$ suffices: a transition moves by one, and the farthest
invariant clause looks two cells to the left. Eight transition rows
enter states having nontrivial clauses. Their 256 Boolean valuations
include 104 satisfying antecedents and 152 inapplicable cases. The other
410 old rows have a trivial destination clause. The initial clause is
vacuous at state~0. This finite check, together with ordinary induction
on computation length, proves the invariant for arbitrary tapes.

The checker separately verifies all 414 ordinary quotient-row equations
and the four exceptional paths. An attempt to read an unspecified local
cell is rejected; the outside tape is genuinely arbitrary. Dense
renumbering supplies 416 further row conjugacy checks. Every one of the
208 working states is graph-reachable from the initial state.

Twelve deliberately damaged controls are rejected. They include an
unredirected incoming edge, a wrong helper, a reversed selector direction,
a changed write bit, corrupted helper branches, an incorrect endpoint
direction, and removal of each needed zero-case context. Normal and
optimized independent-checker receipts are byte-identical. A separate
construction checker enumerates all three-step local histories, adding
all shorter prefixes from the initial state; it independently recovers
the two guards. The theorem is established by the explicit invariant
and path proof, not by a sampled blank computation.

\subsection{Fresh universal primitive replay}
The source and its addresses are unchanged. Nevertheless, we reran the
inherited whole-tape symbolic induction engine on the final 208-state
literal table. This checks the primary candidate directly, in addition
to the local reduction theorem. All 5,084 primitive contracts closed,
using 1,251,865 symbolic inference applications; every one of the 416
literal rows was covered. The smallest constant duration was five, and
all duration coefficients were nonnegative.

The symbolic frame has nine independently arbitrary natural register
values, an arbitrary natural header gap, and an arbitrary nonempty finite
spare-register list. Thirteen affine slots represent the constant, gap,
nine primary values, first spare, and remaining list width. Four
one-block traversal identities are rederived from the actual literal
table; list induction and the unary self-loop rule extend them to
arbitrary block lengths and finite lists. No finite register cutoff is
used in these contracts. The canonical header distance is 56.

Each source increment, positive or zero decrement, branch, continuation,
and halt is checked against its exact whole-tape effect. For nonhalting
primitives the endpoint includes the next encoded source address, head
displacement, all primary values, the spare list, and the blank exteriors.
The 4,000-inference search cap per primitive is a limit on certificate
discovery: every reported case successfully closes with an explicit
finite derivation. It is not a bound on register values or machine runtime.

These contracts are executable mathematical induction evidence. Their
soundness depends on the stated unary/list rules and their implementation;
they are not proof-assistant kernel objects. The new local proof requires
only the much smaller invariant and finite path checker and does not
depend on the universal engine to justify the fold itself.

\subsection{Initializer and independent finite regression}
The parent and new machine execute exactly the same first 48,140
transitions from blank tape. The full trace SHA-256 is
\begin{center}\small\ttfamily
41e9a165abae389c3b182ebbf8d8fd1c260d3d36d4f7ebb6d60e235bac40f19a.
\end{center}
Both arrive in native state~84 with head $-2$, source ordinal 136,
address \texttt{001100110101010}, and 120 one-cells. The logical register
tuple is $[1,0,13,0,0,1,0,0,1]$ and the 43 spare registers all equal one.
The checker verifies the complete finite set of one positions and both
zero exteriors. No changed or deleted row is visited during startup.
This binds initialization precisely; its soundness significance is the
protected suffix argument in Section~\ref{sec:calibration}.

An independent C++ literal executor passed 15,252 primitive cases,
covering three finite assignments for every guarded source branch and
executing 6,778,830 primitive-level literal steps. It also ran a
2,000,000-step blank prefix without halting. The latter observations are
regression evidence only; no finite prefix is offered as a proof of
nonhalting or consistency.

\begin{center}
\small
\begin{tabular}{@{}>{\raggedright\arraybackslash}p{64mm}>{\raggedright\arraybackslash}p{84mm}@{}}
\toprule
Check & Fresh result and scope\\
\midrule
Six-presentation comparison & Two rooted isomorphism classes; identical ordinal source graph\\
New inductive invariant & 104 applicable local Boolean implications; arbitrary tapes by induction\\
New local simulation & 414 ordinary rows and four exhaustive macros, maximum seven steps\\
Dense table and robustness & 208 states, 416 rows, 12 damaged controls rejected; normal/$-O$ agreement\\
Universal runtime replay & 5,084 contracts, all 416 rows, arbitrary natural values and finite lists\\
Exact startup & All 48,140 transitions and complete endpoint tape unchanged\\
Independent finite executor & 15,252 primitive cases; supplementary finite evidence\\
End-to-end kernel certificate & Not produced for U208; inherited ordinary calibration retained\\
\bottomrule
\end{tabular}
\end{center}

\subsection{The alternative 208-state family}
The Compatibility-derived table in \path{alternative/u208.tm} uses the
same deletion with helper~214 and the corresponding old-role labels.
It has 208 working states and 416 rows. Its dense SHA-256 is
\begin{center}\small\ttfamily
9865dcd7bd44deece098c16dcec7b90a96cf111592cbd37e3531b50f7b09276c.
\end{center}
The alternate construction and its history, incoming-census, and local
macro checks pass. We do not claim a separate full universal runtime
rerun for this alternate table; its correctness transfers through its
own exact local theorem from the supplied Compatibility parent.

The two new variants remain nonisomorphic. Let
\[
\begin{split}
p&=\texttt{00000010010000000100110111},\\
q&=\texttt{000000100100000001001101100111111110}.
\end{split}
\]
In the primary table both read-edge graph paths from state~0 end at
native~98. In the alternate table they end at native~98 and native~215,
respectively. A rooted read-edge-preserving bijection would have to
preserve equality of these endpoints, which is impossible. These graph
words need not be realizable as one physical tape history.

\section{Reproduction, artifacts, and further work}
The release is self-contained for its executable checks. Python 3.10 or
later and the standard library suffice for the main replay. No network
access, solver, optimizer, Lean, or Rocq installation is needed.
From the unpacked top-level directory run
\begin{verbatim}
python3 -B reproduce.py
\end{verbatim}
The default command verifies the release hashes, binds the selected
parent and unchanged inference modules to the original Joint archive,
runs that parent's quick construction/source/startup replay in a temporary
directory, reconstructs both 208-state variants, checks the new local
theorem normally and under optimized Python, checks the alternative's
local histories, repeats the six-proposal comparison, verifies exact
startup, and rederives all 5,084 primary U208 contracts. Generated
temporary files do not alter the sealed release. To record a receipt,
pass \texttt{--report} with a path outside the release directory.

\begin{verbatim}
python3 -B reproduce.py --quick
\end{verbatim}
This command explicitly skips the fresh universal runtime derivation.
It remains useful for the new small proof and exact construction, but
must not be reported as a fresh full symbolic replay. The first complete
symbolic run in this preparation environment took about 139 seconds;
wall time is environment-dependent.

The optional independent finite executor requires a C++17 compiler:
\begin{verbatim}
g++ -std=c++17 -O3 code/check_finite_runtime.cpp -o /tmp/check_u208
/tmp/check_u208 machine/u208.native.tm machine/sites.tsv
\end{verbatim}
Rebuild the PDF in a working copy, so generated TeX auxiliary files do not
change the sealed inventory:
\begin{verbatim}
pdflatex -interaction=nonstopmode -halt-on-error report.tex
pdflatex -interaction=nonstopmode -halt-on-error report.tex
\end{verbatim}
PDF bytes may vary with the TeX version or metadata. Literal machine
and source bytes are deterministic.

\begin{center}
\small
\begin{tabular}{@{}>{\raggedright\arraybackslash}p{48mm}>{\raggedright\arraybackslash}p{100mm}@{}}
\toprule
Path & Contents\\
\midrule
\path{report.tex}, \path{report.pdf}, \path{paper/} & Editable report, proofs, comparison, full transition table\\
\path{machine/} & Primary native/dense U208 tables, state map, unchanged source\\
\path{alternative/} & Nonisomorphic Compatibility-derived U208 and its parent/source\\
\path{machine_parent/} & Exact selected U209 table and runtime label correspondence\\
\path{code/}, \path{reproduce.py} & Deterministic builder, independent invariant checker, runtime and startup replay\\
\path{comparison/} & All six table presentations, source snapshots, isomorphisms and fresh quick logs\\
\path{evidence/} & New local, startup, universal, finite, and release replay receipts\\
\path{reference/} & Full selected Joint archive with nested U212/U213/U214/U218/U222 lineage and unary theory; source provenance and reviews\\
\path{research/} & Discovery programs and precisely scoped search receipts\\
\path{licenses/}, \path{NOTICE.txt} & Upstream terms and attribution\\
\path{SHA256SUMS} & Hashes of all other sealed release files\\
\bottomrule
\end{tabular}
\end{center}

\subsection{What the new result suggests}
The effective improvement is a stronger invariant language for state
reuse. An unconditional incoming-bit test misses states entered from
different directions. Here a two-cell old-mode tag and a conditional
old-helper assertion suffice, even though the replacement path revisits
the same physical helper. Complete backward history gives discovery
candidates; small inductive assertions give a clearer proof.

A brief follow-up search on U208 found no further direct alias within
seven-step history contexts and 27 replacement steps, nor a further
shared-row fold across 348 ordered interfaces within the tested
seven-history/25-step limits. These are restricted negative searches,
not a lower bound at 208. Richer invariant combinations, several jointly
changed rows, or high-level arithmetic changes remain possible. The
report does not claim that any of those approaches already gives 207.

The immediate formalization target is compact: encode the five-state
invariant and four finite trace identities in the preferred arithmetic
or proof-assistant framework, then export the generic positive-block
simulation theorem for the exact table. Completing an end-to-end U208
certificate also requires importing or formalizing the retained
numerical-calculus and unary-theory proof transformations. That larger
task should be distinguished from the new finite reduction already
proved here.

\begin{thebibliography}{99}
\bibitem{unary} ProveIt incoming archive, \emph{Unary power-cover theory U}, 29 October 2026.
\href{https://github.com/VladimirReshetnikov/ProveIt/blob/343da82d7384d10ec9fd1c17361fc923b2d1d117/docs/incoming/pa_unary_power_cover.zip}{Pinned source archive}. File \path{pa_unary_power_cover.zip}.
\bibitem{plain} ProveIt incoming archive, \emph{U209 ConPA, Plain}, 29 October 2026.
\href{https://github.com/VladimirReshetnikov/ProveIt/blob/343da82d7384d10ec9fd1c17361fc923b2d1d117/docs/incoming/U209_ConPA_2026-10-29.zip}{Pinned source archive}. File \path{U209_ConPA_2026-10-29.zip}.
\bibitem{compat} ProveIt incoming archive, \emph{U209 ConPA Compatibility}, 29 October 2026.
\href{https://github.com/VladimirReshetnikov/ProveIt/blob/343da82d7384d10ec9fd1c17361fc923b2d1d117/docs/incoming/U209_ConPA_Compatibility_2026-10-29.zip}{Pinned source archive}. File \path{U209_ConPA_Compatibility_2026-10-29.zip}.
\bibitem{folds} ProveIt incoming archive, \emph{U209 ConPA Compatible Folds}, 29 October 2026.
\href{https://github.com/VladimirReshetnikov/ProveIt/blob/343da82d7384d10ec9fd1c17361fc923b2d1d117/docs/incoming/U209_ConPA_Compatible_Folds_2026-10-29.zip}{Pinned source archive}. File \path{U209_ConPA_Compatible_Folds_2026-10-29.zip}.
\bibitem{composed} ProveIt incoming archive, \emph{U209 ConPA Composed}, 29 October 2026.
\href{https://github.com/VladimirReshetnikov/ProveIt/blob/343da82d7384d10ec9fd1c17361fc923b2d1d117/docs/incoming/U209_ConPA_Composed_2026-10-29.zip}{Pinned source archive}. File \path{U209_ConPA_Composed_2026-10-29.zip}.
\bibitem{joint} ProveIt incoming archive, \emph{U209 ConPA Joint, full package}, 29 October 2026.
\href{https://github.com/VladimirReshetnikov/ProveIt/blob/343da82d7384d10ec9fd1c17361fc923b2d1d117/docs/incoming/U209_ConPA_Joint_2026-10-29.zip}{Pinned source archive}. File \path{U209_ConPA_Joint_2026-10-29.zip}.
\bibitem{jointpdf} ProveIt incoming archive, \emph{U209 ConPA Joint, PDF-only presentation}, 29 October 2026.
\href{https://github.com/VladimirReshetnikov/ProveIt/tree/343da82d7384d10ec9fd1c17361fc923b2d1d117/docs/incoming}{Pinned directory}, file \path{U209_ConPA_Joint_2026-10-29 (1).zip}.
\bibitem{k215pa} Elliot Glazer, \emph{K215 halts iff ZFC is inconsistent: Metamath Zero, PA}.
\url{https://github.com/elliotglazer/k215-pa}. Repository documentation inspected 29 October 2026; the large MM0 proof was not recompiled for this report.
\bibitem{k230} Elliot Glazer, \emph{Theory K and the K230 Rocq development}.
\url{https://github.com/elliotglazer/k230-rocq}. In particular, the Theory K section of the README.
\bibitem{m254} The \texttt{jeff87654/busy-beaver} reference repository, \emph{ZF/m254 reference development}.
\url{https://github.com/jeff87654/busy-beaver/tree/zf-254/ZF/m254}. See the original attribution notices preserved in this release.
\end{thebibliography}

\appendix
\clearpage
\section{Complete primary transition table}
\label{app:table}
Working states are $0,\ldots,207$, halt is 208, and the initial state is 0. The head starts at coordinate 0 on the all-zero tape. Each entry lists the written symbol, move, and next state. Both table panels use these same conventions.
\begingroup\small\setlength{\tabcolsep}{5pt}
\begin{longtable}{@{}rll@{\hspace{9mm}}rll@{}}
\toprule
$q$ & Read 0 & Read 1 & $q$ & Read 0 & Read 1\\
\midrule\endfirsthead
\toprule
$q$ & Read 0 & Read 1 & $q$ & Read 0 & Read 1\\
\midrule\endhead
\midrule\multicolumn{6}{r}{\emph{Continued on the next page}}\\\endfoot
\bottomrule\endlastfoot
0 & $1\,R\,92$ & $1\,R\,0$ & 104 & $0\,R\,21$ & $1\,R\,120$ \\
1 & $0\,L\,1$ & $1\,L\,2$ & 105 & $0\,R\,21$ & $0\,L\,10$ \\
2 & $0\,L\,3$ & $1\,L\,2$ & 106 & $0\,R\,61$ & $1\,R\,42$ \\
3 & $0\,L\,5$ & $1\,L\,2$ & 107 & $0\,R\,121$ & $1\,R\,8$ \\
4 & $1\,R\,6$ & $0\,L\,4$ & 108 & $0\,R\,101$ & $1\,R\,122$ \\
5 & $0\,L\,9$ & $1\,L\,2$ & 109 & $0\,L\,207$ & $0\,L\,95$ \\
6 & $0\,R\,6$ & $0\,R\,10$ & 110 & $0\,R\,124$ & $1\,R\,104$ \\
7 & $0\,R\,189$ & $1\,R\,159$ & 111 & $0\,R\,115$ & $1\,L\,72$ \\
8 & $0\,R\,14$ & $1\,R\,13$ & 112 & $0\,R\,194$ & $1\,R\,180$ \\
9 & $0\,L\,15$ & $1\,L\,2$ & 113 & $0\,R\,160$ & $1\,R\,183$ \\
10 & $0\,L\,16$ & $0\,R\,10$ & 114 & $0\,R\,148$ & $1\,R\,146$ \\
11 & $0\,L\,26$ & $1\,R\,128$ & 115 & $0\,R\,11$ & $1\,R\,129$ \\
12 & $0\,R\,19$ & $1\,R\,20$ & 116 & $0\,L\,131$ & $0\,L\,1$ \\
13 & $0\,L\,89$ & $1\,L\,31$ & 117 & $0\,R\,39$ & $1\,R\,72$ \\
14 & $0\,L\,11$ & $0\,R\,22$ & 118 & $0\,R\,61$ & $1\,R\,132$ \\
15 & $0\,L\,23$ & $1\,L\,2$ & 119 & $0\,R\,61$ & $1\,R\,133$ \\
16 & $1\,L\,1$ & $0\,L\,1$ & 120 & $0\,R\,70$ & $0\,L\,16$ \\
17 & $0\,R\,24$ & $0\,R\,32$ & 121 & $0\,R\,14$ & $1\,R\,132$ \\
18 & $0\,R\,98$ & $1\,R\,113$ & 122 & $0\,R\,21$ & $1\,R\,134$ \\
19 & $0\,R\,28$ & $1\,R\,29$ & 123 & $0\,R\,151$ & $1\,R\,188$ \\
20 & $0\,R\,30$ & $1\,R\,31$ & 124 & $0\,R\,137$ & $1\,L\,20$ \\
21 & $0\,R\,24$ & $0\,R\,59$ & 125 & $0\,R\,130$ & $1\,R\,179$ \\
22 & $0\,R\,33$ & $1\,R\,22$ & 126 & $0\,R\,203$ & $1\,R\,156$ \\
23 & $0\,L\,34$ & $1\,L\,2$ & 127 & $0\,R\,75$ & $1\,R\,175$ \\
24 & $0\,L\,35$ & $0\,R\,0$ & 128 & $0\,R\,75$ & $1\,R\,205$ \\
25 & $0\,R\,36$ & $1\,R\,25$ & 129 & $0\,R\,202$ & $1\,R\,157$ \\
26 & $0\,R\,37$ & $1\,R\,38$ & 130 & $0\,R\,12$ & $1\,R\,141$ \\
27 & $0\,R\,39$ & $0\,L\,40$ & 131 & $0\,R\,30$ & $1\,R\,69$ \\
28 & $0\,R\,41$ & $1\,R\,42$ & 132 & $0\,R\,39$ & $1\,L\,124$ \\
29 & $0\,R\,43$ & $1\,R\,44$ & 133 & $0\,R\,153$ & $0\,L\,16$ \\
30 & $0\,R\,45$ & $1\,R\,46$ & 134 & $0\,R\,154$ & $0\,L\,16$ \\
31 & $0\,R\,21$ & $1\,R\,47$ & 135 & $0\,R\,88$ & $1\,R\,136$ \\
32 & $0\,R\,48$ & $1\,R\,32$ & 136 & $0\,R\,196$ & $1\,R\,140$ \\
33 & $0\,R\,32$ & $1\,R\,33$ & 137 & $0\,R\,24$ & $0\,R\,48$ \\
34 & $0\,L\,74$ & $1\,L\,2$ & 138 & $0\,R\,99$ & $1\,R\,37$ \\
35 & $1\,R\,35$ & $1\,L\,1$ & 139 & $0\,R\,171$ & $1\,R\,148$ \\
36 & $0\,R\,50$ & $1\,R\,36$ & 140 & $0\,R\,96$ & $1\,R\,159$ \\
37 & $0\,R\,17$ & $1\,R\,111$ & 141 & $0\,R\,174$ & $1\,R\,165$ \\
38 & $0\,R\,13$ & $1\,R\,52$ & 142 & $0\,R\,12$ & $1\,R\,159$ \\
39 & $0\,L\,11$ & $0\,R\,36$ & 143 & $0\,R\,88$ & $1\,L\,165$ \\
40 & $0\,L\,4$ & $1\,L\,53$ & 144 & $0\,R\,88$ & $1\,R\,158$ \\
41 & $0\,R\,24$ & $0\,R\,25$ & 145 & $0\,L\,152$ & $1\,L\,163$ \\
42 & $0\,R\,13$ & $0\,L\,16$ & 146 & $0\,R\,84$ & $1\,R\,88$ \\
43 & $0\,R\,55$ & $1\,R\,56$ & 147 & $0\,R\,159$ & $1\,R\,12$ \\
44 & $0\,R\,57$ & $1\,R\,58$ & 148 & $0\,R\,83$ & $1\,R\,143$ \\
45 & $0\,R\,24$ & $0\,R\,54$ & 149 & $0\,R\,144$ & $1\,R\,155$ \\
46 & $0\,R\,60$ & $1\,R\,61$ & 150 & $0\,L\,66$ & $1\,L\,141$ \\
47 & $0\,R\,62$ & $0\,L\,16$ & 151 & $0\,R\,150$ & $1\,R\,169$ \\
48 & $0\,R\,25$ & $1\,R\,48$ & 152 & $0\,R\,159$ & $1\,R\,114$ \\
49 & $0\,L\,35$ & $0\,R\,63$ & 153 & $0\,R\,97$ & $1\,R\,97$ \\
50 & $0\,R\,16$ & $1\,R\,50$ & 154 & $0\,R\,49$ & $0\,R\,32$ \\
51 & $0\,R\,147$ & $1\,R\,143$ & 155 & $0\,R\,191$ & $1\,R\,187$ \\
52 & $0\,R\,67$ & $1\,R\,68$ & 156 & $0\,L\,194$ & $1\,R\,194$ \\
53 & $0\,L\,40$ & $1\,L\,53$ & 157 & $0\,L\,186$ & $1\,R\,96$ \\
54 & $0\,R\,22$ & $1\,R\,54$ & 158 & $0\,R\,167$ & $1\,R\,178$ \\
55 & $0\,L\,11$ & $0\,R\,50$ & 159 & $0\,R\,19$ & $1\,R\,170$ \\
56 & $0\,R\,60$ & $0\,L\,16$ & 160 & $0\,R\,164$ & $1\,R\,186$ \\
57 & $0\,R\,69$ & $1\,R\,70$ & 161 & $0\,R\,159$ & $1\,R\,96$ \\
58 & $0\,R\,61$ & $1\,R\,71$ & 162 & $0\,R\,125$ & $1\,R\,164$ \\
59 & $0\,R\,54$ & $1\,R\,59$ & 163 & $0\,R\,97$ & $1\,R\,162$ \\
60 & $0\,R\,49$ & $0\,R\,25$ & 164 & $0\,L\,189$ & $1\,R\,142$ \\
61 & $0\,L\,28$ & $1\,L\,72$ & 165 & $0\,R\,146$ & $1\,R\,171$ \\
62 & $0\,L\,117$ & $1\,L\,30$ & 166 & $0\,R\,96$ & $1\,R\,123$ \\
63 & $0\,R\,73$ & $1\,R\,63$ & 167 & $0\,R\,161$ & $1\,R\,143$ \\
64 & $0\,R\,97$ & $1\,R\,150$ & 168 & $0\,R\,148$ & $1\,R\,87$ \\
65 & $1\,R\,184$ & $1\,L\,53$ & 169 & $0\,R\,161$ & $1\,R\,96$ \\
66 & $0\,R\,144$ & $1\,R\,145$ & 170 & $0\,R\,180$ & $1\,R\,122$ \\
67 & $0\,R\,76$ & $1\,R\,208$ & 171 & $0\,L\,192$ & $1\,R\,60$ \\
68 & $0\,L\,196$ & $1\,R\,31$ & 172 & $0\,R\,145$ & $1\,R\,88$ \\
69 & $0\,R\,41$ & $1\,R\,62$ & 173 & $0\,R\,152$ & $1\,R\,201$ \\
70 & $0\,L\,89$ & $1\,R\,78$ & 174 & $0\,R\,140$ & $1\,R\,84$ \\
71 & $0\,R\,55$ & $1\,L\,79$ & 175 & $0\,R\,174$ & $1\,R\,75$ \\
72 & $0\,R\,80$ & $1\,R\,81$ & 176 & $0\,R\,126$ & $1\,R\,168$ \\
73 & $1\,R\,0$ & $1\,R\,63$ & 177 & $0\,R\,138$ & $1\,R\,116$ \\
74 & $0\,R\,74$ & $1\,R\,82$ & 178 & $0\,R\,97$ & $1\,R\,51$ \\
75 & $0\,L\,145$ & $1\,L\,141$ & 179 & $0\,R\,167$ & $1\,R\,51$ \\
76 & $0\,R\,85$ & $1\,R\,86$ & 180 & $0\,R\,17$ & $1\,L\,131$ \\
77 & $0\,R\,20$ & $1\,R\,38$ & 181 & $0\,R\,143$ & $1\,R\,149$ \\
78 & $0\,R\,89$ & $1\,L\,132$ & 182 & $0\,R\,76$ & $1\,R\,116$ \\
79 & $0\,L\,108$ & $1\,R\,56$ & 183 & $0\,R\,106$ & $1\,R\,124$ \\
80 & $0\,R\,49$ & $0\,R\,33$ & 184 & $1\,L\,40$ & $1\,L\,109$ \\
81 & $0\,R\,91$ & $0\,L\,42$ & 185 & $0\,R\,174$ & $1\,R\,156$ \\
82 & $0\,R\,177$ & $1\,R\,77$ & 186 & $0\,R\,181$ & $1\,R\,182$ \\
83 & $0\,R\,96$ & $1\,R\,12$ & 187 & $0\,R\,66$ & $1\,R\,204$ \\
84 & $0\,R\,88$ & $1\,R\,97$ & 188 & $0\,R\,140$ & $1\,R\,197$ \\
85 & $0\,R\,206$ & $1\,R\,21$ & 189 & $0\,R\,84$ & $1\,R\,190$ \\
86 & $0\,R\,116$ & $1\,R\,100$ & 190 & $0\,R\,157$ & $1\,R\,185$ \\
87 & $0\,R\,173$ & $1\,R\,166$ & 191 & $0\,R\,64$ & $1\,R\,7$ \\
88 & $0\,R\,102$ & $1\,R\,103$ & 192 & $0\,R\,150$ & $1\,R\,127$ \\
89 & $0\,R\,80$ & $1\,R\,79$ & 193 & $0\,R\,18$ & $1\,R\,94$ \\
90 & $0\,R\,101$ & $1\,R\,195$ & 194 & $0\,L\,200$ & $1\,R\,172$ \\
91 & $0\,R\,105$ & $1\,R\,106$ & 195 & $0\,R\,104$ & $1\,R\,96$ \\
92 & $0\,R\,65$ & $0\,R\,0$ & 196 & $0\,R\,90$ & $1\,R\,88$ \\
93 & $0\,R\,27$ & $1\,R\,107$ & 197 & $0\,R\,139$ & $1\,R\,199$ \\
94 & $0\,R\,182$ & $0\,L\,207$ & 198 & $0\,R\,150$ & $1\,R\,112$ \\
95 & $1\,L\,109$ & $1\,L\,95$ & 199 & $0\,R\,201$ & $1\,R\,200$ \\
96 & $0\,R\,19$ & $1\,R\,110$ & 200 & $0\,R\,156$ & $1\,R\,198$ \\
97 & $0\,R\,60$ & $1\,R\,60$ & 201 & $0\,L\,123$ & $1\,R\,167$ \\
98 & $0\,R\,68$ & $1\,R\,135$ & 202 & $0\,R\,196$ & $1\,R\,28$ \\
99 & $0\,R\,93$ & $1\,R\,108$ & 203 & $0\,R\,163$ & $1\,R\,192$ \\
100 & $0\,R\,105$ & $0\,L\,207$ & 204 & $0\,R\,146$ & $1\,R\,176$ \\
101 & $0\,R\,117$ & $1\,R\,118$ & 205 & $0\,R\,161$ & $1\,R\,193$ \\
102 & $0\,R\,71$ & $1\,R\,58$ & 206 & $0\,R\,45$ & $1\,R\,100$ \\
103 & $0\,R\,57$ & $1\,R\,119$ & 207 & $0\,L\,4$ & $0\,L\,207$ \\
\end{longtable}
\endgroup

\end{document}
